\writeSectionFrame{\presentationTitle}{Das zu lösende Problem}
\begin{frame}{Mehrere Objekt anbieten}
	\begin{center}
 		\includegraphics[scale=0.3]{\BILDERPFAD/problem1.png}
 	\end{center}
\end{frame}

\note{
 	 Der Singleton ermöglicht, dass nur ein einziges Objekt angeboten wird.
 	 Aber manchmal möchte man in der Lage sein, mehrere Objekte zu erzeugen,
 	 aber keine beliebige Menge. 
}

\begin{frame}{Eigenes Muster?}
	\begin{center}
 		\includegraphics[scale=0.3]{\BILDERPFAD/problem2.png}
 	\end{center}
\end{frame}

\note{
 	In der Literatur ist es etwas umstritten, ob Multiton wirklich ein eigenes
 	Muster oder nur eine Erweiterung des Singletons ist. Wenn es Ihnen seltsam
 	vorkommt, dass ein Singleton nicht singulär ist, dann nennen Sie ihn besser
 	Multiton.
}